\subsection{Schur Orthogonality Relations}

We keep the notation of the previous subsection. Let \(\lambda\) be an element of \(\widehat{G}\) , and let \(u\) and \(v\) be elements of \(\mathrm{End}_{\mathrm{K}}(\mathrm{V}_{\lambda})\) ; by formula (10), we have

\[
\hat {\tau} (j _ {\lambda} (u) j _ {\lambda} (v)) = d _ {\lambda} \operatorname{Tr} (u v).
\]

Since \(\tau = \hat{\tau} \circ \overline{\mathcal{F}}\) , we deduce from formulas (2) and (15) the relation

\[
d _ {\lambda} ^ {2} | \mathrm{G} | ^ {- 1} \sum_ {g \in \mathrm{G}} \operatorname{Tr} (u \pi_ {\lambda} (g)) \operatorname{Tr} (v \pi_ {\lambda} (g ^ {- 1})) = d _ {\lambda} \operatorname{Tr} (u v).
\]

We observed previously that \(d_{\lambda}.1 \neq 0\) in the field K; it follows that

\[
| \mathrm{G} | ^ {- 1} \sum_ {g \in \mathrm{G}} \operatorname{Tr} (u \pi_ {\lambda} (g)) \operatorname{Tr} (v \pi_ {\lambda} (g ^ {- 1})) = d _ {\lambda} ^ {- 1} \operatorname{Tr} (u v).\tag{20}
\]

Let us apply this relation to the case when u and v have rank \(\leqslant 1\) ; we obtain

\[
| \mathrm{G} | ^ {- 1} \sum_ {g \in \mathrm{G}} \left\langle x ^ {*}, \pi_ {\lambda} (g) x \right\rangle \left\langle y ^ {*}, \pi_ {\lambda} \left(g ^ {- 1}\right) y \right\rangle = d _ {\lambda} ^ {- 1} \left\langle x ^ {*}, y \right\rangle \left\langle y ^ {*}, x \right\rangle\tag{21}
\]

for \(x, y\) in \(\mathrm{V}_{\lambda}\) and \(x^{*}, y^{*}\) in the dual \(\mathrm{V}_{\lambda}^{*}\) of \(\mathrm{V}_{\lambda}\) .

For every \(\lambda \in \widehat{\mathbf{G}}\) , let \((e_{\lambda,j})_{1\leqslant j\leqslant d_{\lambda}}\) be a basis of \(\mathrm{V}_{\lambda}\) ; denote by \((\pi_{ij}^{\lambda}(g))\) the matrix of the endomorphism \(\pi_{\lambda}(g)\) of \(\mathrm{V}_{\lambda}\) with respect to this basis. If we denote by \((e_{\lambda,i}^{*})_{1\leqslant i\leqslant d_{\lambda}}\) the basis of \(\mathrm{V}_{\lambda}^{*}\) dual to \((e_{\lambda,j})\) , then we have \(\pi_{ij}^{\lambda}(g) = \langle e_{\lambda,i}^{*}, \pi_{\lambda}(g)e_{\lambda,j}\rangle\) , and therefore

\[
| \mathrm{G} | ^ {- 1} \sum_ {g \in \mathrm{G}} \pi_ {i j} ^ {\lambda} (g) \pi_ {k \ell} ^ {\lambda} (g ^ {- 1}) = d _ {\lambda} ^ {- 1} \delta_ {i \ell} \delta_ {j k}.\tag{22}
\]

Now let \(\lambda\) and \(\mu\) be two distinct elements of \(\widehat{\mathbf{G}}\) , and let \(u \in \mathrm{End}_{\mathrm{K}}(\mathrm{V}_{\lambda})\) and \(v \in \mathrm{End}_{\mathrm{K}}(\mathrm{V}_{\mu})\) . Again by relation (15), we have

\[
\sum_ {g \in \mathrm{G}} \operatorname{Tr} (u \pi_ {\lambda} (g)) \operatorname{Tr} (v \pi_ {\mu} (g ^ {- 1})) = 0.\tag{23}
\]

As above, we deduce that

\[
\sum_ {g \in \mathrm{G}} \langle x ^ {*}, \pi_ {\lambda} (g) x \rangle \langle y ^ {*}, \pi_ {\mu} (g ^ {- 1}) y \rangle = 0\tag{24}
\]

for \(x\in \mathrm{V}_{\lambda},x^{*}\in \mathrm{V}_{\lambda}^{*},y\in \mathrm{V}_{\mu}\) , and \(y^{*}\in \mathrm{V}_{\mu}^{*}\) . We also have

\[
\sum_ {g \in \mathrm{G}} \pi_ {i j} ^ {\lambda} (g) \pi_ {k \ell} ^ {\mu} (g ^ {- 1}) = 0\tag{25}
\]

for \(i,j\) in \([1,d_{\lambda}]\) and \(k,\ell\) in \([1,d_{\mu}]\) .

Relations (20) through (25) are known as the Schur orthogonality relations.

Remark. --- We identify the algebra \(\operatorname{End}_{\mathrm{K}}(\mathrm{V}_{\lambda})\) with the matrix algebra \(\mathbf{M}_{d_{\lambda}}(\mathrm{K})\) via the basis \((e_{\lambda,j})\) of \(V_{\lambda}\) . The mapping \(\overline{F}^{-1}\) is an isomorphism from the algebra \(\prod_{\lambda}\mathbf{M}_{d_{\lambda}}(\mathrm{K})\) to the algebra K{[}G{]}. For \(\mu\in\widehat{G}\) , denote by \(E_{ij}^{\mu}\) the element of \(\prod_{\lambda}\mathbf{M}_{d_{\lambda}}(\mathrm{K})\) whose component of index \(\mu\) is the matrix unit \(E_{ij}\) of \(\mathbf{M}_{d_{\mu}}(\mathrm{K})\) (II, §10, No.~3, p.~341) and whose other components are zero; set \(u_{ij}^{\mu}=\overline{\mathcal{F}}^{-1}(\mathrm{E}_{ij}^{\mu})\) . The family of elements \(u_{ij}^{\lambda}\) , for \(\lambda\in\widehat{G}\) , \(1\leqslant i\leqslant d_{\lambda}\) , \(1\leqslant j\leqslant d_{\lambda}\) , is a basis of the algebra K{[}G{]}; the multiplication table is

\[
u _ {i j} ^ {\lambda} u _ {k \ell} ^ {\mu} = \delta_ {\lambda \mu} \delta_ {j k} u _ {i \ell} ^ {\lambda}.\tag{26}
\]

Moreover, by formula (15), we have

\[
u _ {i j} ^ {\lambda} = | \mathrm{G} | ^ {- 1} d _ {\lambda} \sum_ {g \in \mathrm{G}} \pi_ {j i} ^ {\lambda} (g ^ {- 1}) g.\tag{27}
\]

